% Extensión aditiva. Canon leído: CURRENT.json rev. 2026-07-22.2.
% Autores del proyecto: Oumar Haidara Fall y Rubén Ramos Balsa.

\chapter{Rectas de magnitud y realización dimensional del núcleo HMT}
\label{chap:funtor-dimensional-hmt}

La transición desde HMT hacia Mecánica Dimensional no consiste en adjuntar
un símbolo de unidad a un número real. El núcleo aritmético produce
caracteres, cociclos, razones y órdenes de una carta decimal. Una magnitud
física, en cambio, es un elemento de una recta provista de tipo dimensional.
La distinción es estructural: un cambio de unidades modifica las coordenadas
de una magnitud, pero no la magnitud misma. Esta sección construye, por una
parte, la retícula ambiental requerida para realizar independientemente
masa, longitud, duración y carga, y determina que su rango mínimo es cuatro.
Por otra parte, construye un funtor monoidal de trayectorias con firma que
toma valores en una porción de esa retícula. La minimalidad de la retícula no
se atribuye al funtor.

La construcción tiene dos consecuencias. En primer lugar, permite realizar
en una retícula común la acción, la duración, la longitud, la carga y la masa
sin introducir el Sistema Internacional en las fórmulas generadoras; la
asignación de carga a estados requiere, además, un carácter entero
específico. En segundo lugar, demuestra exactamente qué parte de una escala
absoluta no puede proceder de escalares adimensionales. Esta imposibilidad no
afecta a las razones, a los caracteres de masa ni a las identidades
constitutivas; fija el tipo matemático correcto de sus salidas.
La distinción entre dimensión, coordenada numérica y elección de unidad es la
de la teoría clásica del análisis dimensional \cite{Buckingham1914,Bridgman1931};
la estructura de rectas que sigue hace explícita esa distinción en el dominio
HMT\@.

\section{Retícula de dimensiones y rectas orientadas}

Sea
\[
 \Lambda_{\rm D}
 =\mathbb Z\langle \mathbf M,\mathbf L,\mathbf T,\mathbf Q\rangle
 \cong\mathbb Z^4
\]
la retícula de dimensiones, escrita en términos de masa, longitud, duración
y carga. Para
\(d=(d_M,d_L,d_T,d_Q)\in\Lambda_{\rm D}\), denotamos por
\(\mathcal L_d\) una recta real orientada y por
\(\mathcal L_d^+\) su semirrecta positiva. El producto tensorial realiza la
suma de dimensiones,
\[
 \mathcal L_d\otimes\mathcal L_e\simeq\mathcal L_{d+e},
 \qquad
 \mathcal L_d^\vee\simeq\mathcal L_{-d}.
\]
La colección de estas rectas, con isomorfismos positivos y producto
tensorial, forma un grupoide de Picard simétrico. Un escalar adimensional es
un elemento de \(\mathcal L_0\simeq\mathbb R\); una magnitud de dimensión
\(d\) es un elemento de \(\mathcal L_d\).

Introducimos cuatro dimensiones adaptadas a HMT:
\[
 \mathbf S=\mathbf M+2\mathbf L-\mathbf T,
 \qquad
 \mathbf C=\mathbf L-\mathbf T,
 \qquad
 \mathbf T=\mathbf T,
 \qquad
 \mathbf Q=\mathbf Q.
\]
Representan, respectivamente, acción, velocidad, duración y carga. En la
base \((\mathbf M,\mathbf L,\mathbf T,\mathbf Q)\), sus columnas forman la
matriz
\[
 B_{\rm D}=
 \begin{pmatrix}
 1&0&0&0\\
 2&1&0&0\\
 -1&-1&1&0\\
 0&0&0&1
 \end{pmatrix}.
\]

\begin{theorem}[Base dimensional adaptada a HMT]
\label{thm:base-dimensional-hmt}
Las dimensiones
\[
 \mathbf S,\qquad\mathbf C,\qquad\mathbf T,\qquad\mathbf Q
\]
constituyen una base unimodular de
\(\Lambda_{\rm D}\). En particular, toda dimensión mecánica,
electromagnética o gravitatoria generada por
\((\mathbf M,\mathbf L,\mathbf T,\mathbf Q)\) posee una expresión única con
exponentes enteros en las cuatro rectas
\(\mathcal L_S,\mathcal L_C,\mathcal L_T,\mathcal L_Q\).
\end{theorem}

\begin{proof}
La matriz \(B_{\rm D}\) es triangular por bloques y
\(\det B_{\rm D}=1\). Por ello pertenece a
\(\operatorname{GL}_4(\mathbb Z)\). Su inversa integral proporciona
\[
 \mathbf L=\mathbf C+\mathbf T,
 \qquad
 \mathbf M=\mathbf S-2\mathbf C-\mathbf T.
\]
Las cuatro dimensiones originales se recuperan, por tanto, de forma única.
\end{proof}

Entre las identidades que utilizaremos se encuentran
\begin{align}
 \mathbf E&=\mathbf S-\mathbf T,
 &\mathbf Z&=\mathbf S-2\mathbf Q,\label{eq:dim-EZ}\\
 \boldsymbol\mu&=\mathbf S-\mathbf C-2\mathbf Q,
 &\boldsymbol\varepsilon&=-\mathbf S-\mathbf C+2\mathbf Q,
 \label{eq:dim-mueps}\\
 \mathbf G&=-\mathbf S+5\mathbf C+2\mathbf T,
 &\boldsymbol\kappa&=-\mathbf S+\mathbf C+2\mathbf T,
 \label{eq:dim-Gkappa}\\
 \boldsymbol\Lambda&=-2\mathbf C-2\mathbf T.
 \label{eq:dim-Lambda}
\end{align}
Aquí \(\mathbf E\) es energía, \(\mathbf Z\) impedancia,
\(\boldsymbol\mu\) permeabilidad, \(\boldsymbol\varepsilon\)
permitividad, \(\mathbf G\) la dimensión de la constante gravitatoria,
\(\boldsymbol\kappa=\mathbf G-4\mathbf C\), y
\(\boldsymbol\Lambda=-2\mathbf L\).

\section{Grupo de reescala y teorema de imposibilidad}

El grupo de cambios de base de unidades es
\[
 \mathcal G_{\rm D}
 :=\operatorname{Hom}(\Lambda_{\rm D},\mathbb R_{>0})
 \cong(\mathbb R_{>0})^4.
\]
Para \(g\in\mathcal G_{\rm D}\) y \(d\in\Lambda_{\rm D}\), escribimos
\(g^d\) para el carácter correspondiente. Si \(u_d\) es una base positiva
de \(\mathcal L_d\), el cambio de carta
\(u_d\mapsto g^d u_d\) transforma la coordenada de una magnitud fija según
\[
 x\,u_d=(g^{-d}x)(g^du_d).
\]
La magnitud no cambia; cambia su representación numérica.

\begin{theorem}[Imposibilidad de una escala absoluta desde datos
adimensionales]
\label{thm:no-go-dimensional-rango-cuatro}
Sea \(X\) un espacio sobre el que \(\mathcal G_{\rm D}\) actúa
trivialmente y sea \(d\ne0\). Toda aplicación
\(F:X\to\mathcal L_d\) que sea natural respecto de todos los cambios de
unidades y que no reciba una base dimensional adicional es idénticamente
nula. En cambio, todo cociente de dos secciones no nulas de una misma recta,
y todo producto cuya dimensión total sea cero, es invariante por
\(\mathcal G_{\rm D}\).
\end{theorem}

\begin{proof}
La naturalidad y la acción trivial sobre \(X\) implican
\[
 F(x)=F(gx)=g^dF(x)
 \qquad(g\in\mathcal G_{\rm D}).
\]
Como \(d\ne0\), existe \(g\) con \(g^d\ne1\); de aquí
\((g^d-1)F(x)=0\), luego \(F(x)=0\). Para dos secciones de dimensión
\(d\), los factores \(g^d\) se cancelan en el cociente. La misma cancelación
ocurre en cualquier monomio de dimensión total cero.
\end{proof}

El teorema no establece que una teoría discreta sea incapaz de construir una
escala interna. Establece que esa escala debe aparecer como una sección de
una recta o, equivalentemente, que la teoría debe extender su dominio por el
torsor de marcos dimensionales. Es precisamente esta extensión la que define
la Mecánica Dimensional.

\begin{corollary}[Rango mínimo]
Todo sistema de rectas que permita realizar independientemente masa,
longitud, duración y carga posee rango de escala al menos cuatro. La base
\((\mathbf S,\mathbf C,\mathbf T,\mathbf Q)\) alcanza la cota; por tanto es
mínima.
\end{corollary}

\begin{proof}
Las cuatro dimensiones originales son linealmente independientes en
\(\Lambda_{\rm D}\). Tres generadores no pueden abarcar una retícula de
rango cuatro. El \cref{thm:base-dimensional-hmt} da cuatro generadores que
sí la abarcan.
\end{proof}

\section{Espacio dimensional interno: grados, cartas y operadores de transducción HMT–MD}

Un marco interno HMT--MD es una cuaterna de bases positivas
\[
 \mathfrak u=(u_S,u_C,u_T,u_Q)
 \in\operatorname{Fr}_{\rm D},
\]
donde \(\operatorname{Fr}_{\rm D}\) es un torsor principal de
\(\mathcal G_{\rm D}\). Las bases derivadas se definen tensorialmente:
\begin{align*}
 u_L&=u_Cu_T,&
 u_M&=u_Su_C^{-2}u_T^{-1},&
 u_E&=u_Su_T^{-1},\\
 u_Z&=u_Su_Q^{-2},&
 u_\mu&=u_Su_C^{-1}u_Q^{-2},&
 u_\varepsilon&=u_S^{-1}u_C^{-1}u_Q^2,\\
 u_G&=u_S^{-1}u_C^5u_T^2,&
 u_\kappa&=u_S^{-1}u_Cu_T^2,&
 u_\Lambda&=u_C^{-2}u_T^{-2}.
\end{align*}
La yuxtaposición denota aquí producto tensorial de bases. Una carta externa
es otro marco \(\mathfrak u_{\rm ext}\) junto con un isomorfismo positivo
entre ambos sistemas de rectas. Nada obliga a que esa carta sea SI; la
elección de SI pertenece a la publicación metrológica de una magnitud.

\subsection{Recta de acción}

El lector determinantal local parte de
\[
 G_H=\operatorname{coker}(H_4),
 \qquad |G_H|=16,
 \qquad \Sigma_H=\varphi\alpha_{\rm HMT}^{16},
\]
y demuestra
\[
 10^{-34}<\Sigma_H<10^{-33},
 \qquad \operatorname{ord}_{10}(\Sigma_H)=-34.
\]
La realización regional posterior al retorno proporciona la coordenada
angular adimensional
\[
 \etaRet
 =\frac{\Cdeg}{2}-\alpha_{\angle,{\rm deg}}
 =1.0545718169150390611336905847242997\ldots.
\]
Definimos la sección interna de acción reducida por
\begin{equation}
 \hbar_{\rm int}
 :=\etaRet\,10^{-34}u_S,
 \qquad
 h_{\rm int}:=2\pi\hbar_{\rm int}.
 \label{eq:accion-interna}
\end{equation}
El exponente es una salida del lector; \(u_S\) es la base de una recta.
La composición es covariante bajo \(\mathcal G_{\rm D}\) y no contiene una
identificación con \(\mathrm{J\,s}\).

\subsection{Periodicidad de aristas de ruta y velocidad constitutiva resultante}

Sea \(\mathsf P_{\rm TPK}\) la categoría libre de trayectorias del grafo de
estados enriquecidos. La graduación por longitud
\[
 \tau(\gamma)=|\gamma|,
 \qquad
 \tau(\gamma\eta)=\tau(\gamma)+\tau(\eta)
\]
define el reloj combinatorio del protocolo. Su realización es
\[
 T_{\rm int}(\gamma)=\tau(\gamma)u_T.
\]
Esta duración cuenta transiciones del grafo fijado. Una subdivisión del grafo
define otro reloj y debe acompañarse de su morfismo de refinamiento.

Los canales \(q_\pm\) generan las funciones simétrica y antisimétrica
\[
 \mathcal E=
 \log\!\bigl[(1-q_+^{90})(1-q_-^{90})\bigr]
 -\log\!\bigl[(1-q_+^{120})(1-q_-^{120})\bigr],
\]
\[
 \mathcal B=
 \log\!\frac{1-q_-^{90}}{1-q_+^{90}}
 -\log\!\frac{1-q_-^{120}}{1-q_+^{120}}.
\]
Éstas determinan
los coeficientes positivos
\[
 \widehat c=e^{-\mathcal E},
 \qquad \widehat Z=e^{\mathcal B}.
\]
Sus realizaciones internas son
\begin{equation}
 c_{\rm int}=\widehat c\,u_C,
 \qquad Z_{\rm int}=\widehat Z\,u_Z.
 \label{eq:cZ-internos}
\end{equation}
Para un canal constante a lo largo de una trayectoria se obtiene la longitud
\begin{equation}
 L_{\rm int}(\gamma)
 =\widehat c\,|\gamma|\,u_L.
 \label{eq:longitud-ruta}
\end{equation}
Si el canal depende de la arista, la fórmula covariante es la suma
\(\sum_{e\in\gamma}\widehat c(e)u_L\); su aditividad por concatenación se
conserva.

\subsection{Escala elemental de carga y carácter condicional}

Puesto que \(u_Z=u_Su_Q^{-2}\), el cociente
\(h_{\rm int}/Z_{\rm int}\) pertenece a \(\mathcal L_Q^{\otimes2}\).
La orientación positiva de las rectas selecciona una única raíz positiva.
Definimos
\begin{equation}
 e_{\rm int}
 :=\left(\frac{2\alpha_{\rm HMT}h_{\rm int}}
 {Z_{\rm int}}\right)^{1/2}\in\mathcal L_Q^+.
 \label{eq:carga-interna}
\end{equation}

\begin{proposition}[Clausura interna de la relación de estructura fina]
La sección \eqref{eq:carga-interna} es la única sección positiva que
satisface
\[
 \alpha_{\rm HMT}
 =\frac{Z_{\rm int}e_{\rm int}^{\,2}}{2h_{\rm int}}.
\]
La igualdad es invariante bajo el grupo completo de reescala.
\end{proposition}

\begin{proof}
La existencia resulta de la positividad de
\(\alpha_{\rm HMT},h_{\rm int},Z_{\rm int}\). Una recta real orientada
posee una sola raíz positiva de una sección positiva de su cuadrado. La
dimensión del miembro derecho es
\(\mathbf Z+2\mathbf Q-\mathbf S=0\), de modo que todos los factores de
reescala se cancelan.
\end{proof}

Esta proposición no cuenta como una segunda determinación independiente de
\(\alpha\): realiza, dentro del sistema de rectas, la compatibilidad
constitutiva entre acción, impedancia y una escala elemental de carga. No
construye una carga sobre los estados. Para ello se necesita además un
carácter entero
\[
 n_Q:\operatorname{Ob}(\mathsf P_{\rm TPK}^{\rm sig})\longrightarrow
 \mathbb Z;
 \qquad
 Q(x)=n_Q(x)e_{\rm int}.
\]
La existencia y la selección de \(n_Q\) son hipótesis adicionales. En
particular, \(n_Q\) no es una componente del funtor de trayectorias que se
construye a continuación. La escala elemental y la clasificación de estados
son mapas distintos.

\section{Funtor monoidal de trayectorias con firma}

Sea \(\mathsf P_{\rm TPK}^{\rm sig}\) la categoría de trayectorias
enriquecidas decoradas por una coordenada de firma
\(v_x\in\mathbb Z^5\) en cada objeto. Para
\(\gamma:x\to y\), ponemos
\[
 \Delta v(\gamma)=v_y-v_x.
\]
La composición satisface
\(\Delta v(\gamma\eta)=\Delta v(\gamma)+\Delta v(\eta)\).
Sea
\[
 \chi_{\rm mass}(v)=\exp\langle\lambda_{\rm mass},v\rangle
\]
el carácter positivo de masa. Sobre el cociente de rutas fijamos la ley de
acción determinada por cinco datos: la traza normalizada de las emisiones
distintas, un cociclo positivo y aditivo, el defecto de cambio
\(\Delta_4P_3\), la capacidad residual \(\lambda=1-\delta\) y la suma local
de los costes de arista. En esta ley,
\[
 \mathcal A_*(\gamma)
 =N_{\rm giro}(\gamma)+\lambda_*N_{\rm hoja}(\gamma),
 \qquad
 \lambda_*=1-\frac5{468}-\frac3{11}\Delta_4.
\]
El argumento siguiente sólo requiere la aditividad; también se aplica al
funcional \(I_\lambda\) con cualquier \(\lambda>0\).

Consideremos el monoide
\[
 \mathfrak M_{\rm D}
 =\mathbb N\times\mathbb R_{\ge0}\times
  \mathbb R_{\ge0}\times\mathbb R_{>0}
\]
con producto
\[
 (t,\ell,s,r)\odot(t',\ell',s',r')
 =(t+t',\ell+\ell',s+s',rr').
\]

\begin{theorem}[Funtor monoidal de trayectorias con firma]
\label{thm:funtor-dimensional-minimo}
La aplicación
\[
 \mathfrak D_{\rm HMT}(\gamma)
 =\bigl(
 |\gamma|,
 \widehat c|\gamma|,
 \mathcal A_*(\gamma),
 \chi_{\rm mass}(\Delta v(\gamma))
 \bigr)
\]
define un funtor monoidal
\[
 \mathfrak D_{\rm HMT}:
 \mathsf P_{\rm TPK}^{\rm sig}\longrightarrow\mathfrak M_{\rm D}.
\]
Después de elegir un marco interno, sus tres componentes aditivas se realizan
como
\[
 T_{\rm int}(\gamma)=|\gamma|u_T,
 \quad
 L_{\rm int}(\gamma)=\widehat c|\gamma|u_L,
 \quad
 S_{\rm int}(\gamma)=\mathcal A_*(\gamma)\hbar_{\rm int},
\]
y la cuarta transporta razones adimensionales de masa. El funtor no contiene
una coordenada de carga ni genera la retícula dimensional ambiental.
\end{theorem}

\begin{proof}
La longitud de una trayectoria, el número de giros y el número de cambios de
hoja son aditivos bajo concatenación compatible. Por hipótesis,
\(\mathcal A_*\) es aditiva. La diferencia de firmas es aditiva y el carácter
de masa transforma sumas en productos. Cada componente respeta, por tanto,
la ley \(\odot\), y la ruta vacía se envía a
\((0,0,0,1)\). La realización lineal conserva las sumas porque multiplica
por secciones fijas de las rectas correspondientes.
\end{proof}

\begin{remark}[Retícula ambiental y alcance del funtor]
El corolario de rango mínimo se refiere a un sistema capaz de variar
independientemente masa, longitud, duración y carga. El funtor anterior
realiza duración, longitud y acción, y transporta una razón de masa; no usa
\(u_Q\) ni demuestra una cota de rango. La posterior asignación de cargas
permanece condicionada a la existencia de \(n_Q\). En la retícula dimensional,
los tipos de sus componentes son
\[
 \mathbf T,\qquad \mathbf L=\mathbf C+\mathbf T,\qquad
 \mathbf S,\qquad \mathbf0,
\]
por lo que su imagen dimensional está contenida en
\(\langle\mathbf S,\mathbf C,\mathbf T\rangle_{\mathbb Z}\), de rango tres.
El reescalado \(u_Q\mapsto\lambda u_Q\), con las otras bases fijas, deja
invariantes todas las componentes del funtor; éste no realiza fielmente la
cuarta dirección ambiental.
\end{remark}

Si a cada objeto se le asigna la sección electrónica bidireccional como
origen de la carta de masa,
\[
 m(x)=m_e^{\beta}\chi_{\rm mass}(v_x-v_e),
\]
entonces todo morfismo \(\gamma:x\to y\) satisface exactamente
\[
 \frac{m(y)}{m(x)}
 =\chi_{\rm mass}(\Delta v(\gamma)).
\]
El teorema cierra la propagación dimensional una vez dada la firma. No cambia
el estatuto del mapa anterior que asocia una especie a una trayectoria y a
su firma.

\section{Etapa basal y clausura bidireccional de la escala electrónica}

El invariante electrónico vigente es el escalar positivo
\[
 \mathfrak e_0=
 \frac{\sqrt3}{4}
 \left(e^{\varphi/\pi^2}-22\alpha_{\rm HMT}^3\right)
 (1+15\Delta_4).
\]
Su valor numérico es
\(0.510998922488066\ldots\). Ésta es la etapa basal de la genealogía, no su
salida rectora. En el marco dimensional interno define una sección de energía
y una sección de masa relacionadas por la velocidad:
\begin{equation}
 E_e^{\rm int}=\mathfrak e_0u_E,
 \qquad
 m_e^{\rm int}=\frac{E_e^{\rm int}}{c_{\rm int}^{2}}
 =\mathfrak e_0\widehat c^{-2}u_M.
 \label{eq:electron-line-valued}
\end{equation}
La ruta central y sus dos lectores bidireccionales determinan, sin consultar
la masa externa,
\[
 \kappa_e=80-54\alpha_{\rm HMT}+6\Delta_4,
 \qquad R_{\rm act}=\frac{H_5}{\etaRet}.
\]
La sección electrónica vigente es, por tanto,
\begin{equation}
 \begin{aligned}
 E_e^{\beta}&=E_e^{(0)}R_{\rm act}^{\kappa_e},
 &m_e^{\beta}&=\frac{E_e^{\beta}}{c_{\rm int}^{2}},\\
 E_e^\beta
 &=0.5109989506900008086082571648\ldots\ {\rm MeV},
 &u_E&=1\,{\rm MeV}.
 \end{aligned}
 \label{eq:electron-line-valued-beta}
\end{equation}
Así, la fórmula basal conserva su función derivativa y la clausura
bidireccional gobierna toda propagación posterior por caracteres y torres.
En una carta de unidades naturales en la que la coordenada de
\(c_{\rm int}\) se fija en uno, ambas coordenadas numéricas coinciden; sus
rectas no se identifican silenciosamente.

\begin{proposition}[Separación tipada de la base energética y el cuanto de reloj]
\label{prop:no-identificacion-base-reloj-rev8}
Escribamos \(t_0=\tau_0u_T\), con \(\tau_0>0\), y definamos
\[
 \varepsilon_{\rm clk}
 :=\hbar_{\rm int}\frac{2\pi}{108t_0},
 \qquad
 \kappa_{\rm clk}(\tau_0)
 :=\frac{2\pi\etaRet\,10^{-34}}{108\tau_0}.
\]
Entonces
\[
 \varepsilon_{\rm clk}=\kappa_{\rm clk}(\tau_0)u_E,
 \qquad
 m_0:=\frac{\varepsilon_{\rm clk}}{c_{\rm int}^2}
 =\kappa_{\rm clk}(\tau_0)\widehat c^{-2}u_M,
\]
y la sección electrónica vigente satisface
\[
 E_e^{\rm int}
 =\frac{\mathfrak e_0}{\kappa_{\rm clk}(\tau_0)}
  \varepsilon_{\rm clk}.
\]
Por tanto, \(E_e^{\rm int}=\mathfrak e_0\varepsilon_{\rm clk}\) si y sólo si
\(\kappa_{\rm clk}(\tau_0)=1\). Adoptar \(\varepsilon_{\rm clk}\) como base
manteniendo la misma coordenada \(\mathfrak e_0\) es una nueva elección de
marco; no constituye una corrección del invariante electrónico ni se sigue
de la periodicidad de orden 108. En particular, la normalización usada más
adelante, \(t_0=u_T\), fija \(\tau_0=1\), no
\(\kappa_{\rm clk}=1\).
\end{proposition}

\begin{proof}
Se sustituyen
\(\hbar_{\rm int}=\etaRet\,10^{-34}u_S\),
\(t_0=\tau_0u_T\), \(u_E=u_Su_T^{-1}\) y
\(c_{\rm int}=\widehat c,u_C\). La última identidad resulta al comparar
con \cref{eq:electron-line-valued}. Todas las igualdades se producen entre
secciones de la misma recta; ninguna consulta una unidad o masa externa.
\end{proof}

\begin{proposition}[Propagación de la escala de masa]
Para cada firma declarada \(v\in\mathbb Z^5\),
\[
 m_v^{\rm int}:=m_e^{\rm int}\chi_{\rm mass}(v)
 \]
es una sección positiva de \(\mathcal L_M\), y
\[
 \frac{m_{v+w}^{\rm int}}{m_v^{\rm int}}
 =\chi_{\rm mass}(w).
\]
La construcción es covariante bajo reescala y no requiere una unidad externa
de energía o masa.
\end{proposition}

\begin{proof}
El carácter es adimensional y positivo, por lo que no modifica la recta de
\(m_e^{\rm int}\). La identidad de cocientes es la multiplicatividad de
\(\chi_{\rm mass}\). Bajo un cambio de marco, numerador y denominador
adquieren el mismo peso \(\mathbf M\), que se cancela.
\end{proof}

Una carta metrológica posterior puede enviar \(u_E\) a una base de energía y
\(u_M\) a la base coherente obtenida al dividir por la velocidad patrón al
cuadrado. Sólo esa carta autoriza expresar la sección mediante un nombre de
unidad externo.

\section{Constitutivas como isomorfismos entre rectas}

Las funciones internas satisfacen
\[
 \widehat\mu=e^{\mathcal B+\mathcal E},
 \qquad
 \widehat\varepsilon=e^{-\mathcal B+\mathcal E},
 \qquad
 \widehat Z=e^{\mathcal B},
 \qquad
 \widehat c=e^{-\mathcal E}.
\]
Definimos las secciones
\[
 \mu_{\rm int}=\widehat\mu u_\mu,
 \qquad
 \varepsilon_{\rm int}=\widehat\varepsilon u_\varepsilon.
\]

\begin{theorem}[Constitutivas lineales covariantes]
En el grupoide de rectas orientadas se cumplen las identidades exactas
\[
 \sqrt{\frac{\mu_{\rm int}}{\varepsilon_{\rm int}}}
 =Z_{\rm int},
 \qquad
 \frac1{\sqrt{\mu_{\rm int}\varepsilon_{\rm int}}}
 =c_{\rm int}.
\]
Las cuatro secciones dependen de dos caracteres adimensionales y de las dos
rectas \(\mathcal L_Z,\mathcal L_C\); no constituyen cuatro escalas
independientes.
\end{theorem}

\begin{proof}
En coeficientes,
\(\widehat\mu/\widehat\varepsilon=\widehat Z^2\) y
\(\widehat\mu\widehat\varepsilon=\widehat c^{-2}\). En dimensiones,
\[
 \boldsymbol\mu-\boldsymbol\varepsilon=2\mathbf Z,
 \qquad
 \boldsymbol\mu+\boldsymbol\varepsilon=-2\mathbf C.
\]
Las raíces positivas son, por ello, secciones de las rectas correctas y las
identidades son covariantes.
\end{proof}

\section{Escala de Cartan y homogeneidad de la acción gravitatoria}

La duración elemental \(t_0=u_T\) y la velocidad constitutiva determinan la
longitud de una transición,
\[
 \ell_0=c_{\rm int}t_0=\widehat c\,u_L.
\]
Esta sección permite homogeneizar la conexión afín de Cartan:
\[
 \mathcal A_{\ell_0}
 =\begin{pmatrix}
 \omega^I{}_J&\ell_0^{-1}e^I\\
 0&0
 \end{pmatrix},
 \qquad
 \mathcal F_{\ell_0}
 =\begin{pmatrix}
 F^I{}_J&\ell_0^{-1}T_{\rm tor}^I\\
 0&0
 \end{pmatrix}.
\]
El cociente \(e/\ell_0\) es adimensional; curvatura y torsión permanecen en
bloques distintos del mismo transporte.

La dimensión de \(G\) no añade un quinto generador, pues
\(\mathbf G=-\mathbf S+5\mathbf C+2\mathbf T\). Para todo parámetro
adimensional positivo \(\theta_G\), definimos la familia covariante
\begin{equation}
 G_{\rm int}(\theta_G)
 :=\theta_G\frac{c_{\rm int}^{5}t_0^2}{\hbar_{\rm int}}
 \in\mathcal L_G^+.
 \label{eq:G-interno}
\end{equation}
La elección \(\theta_G=1\) es la normalización interna de Planck asociada a
una transición elemental. Satisface por construcción
\[
 \ell_0^2=\frac{\hbar_{\rm int}G_{\rm int}(1)}
 {c_{\rm int}^{3}}.
\]
Esta normalización de referencia no es el selector circular ni el selector
holonómico de la sección física de \(G\).
El parámetro \(\theta_G\) hace visible la libertad que no fijan las
identidades dimensionales. Una dinámica TPK que produzca un carácter
gravitatorio deberá determinarlo antes de cualquier carta externa.

Sea
\[
 \kappa_{\rm int}=\frac{8\pi G_{\rm int}}{c_{\rm int}^{4}}.
\]
En el funcional siguiente, \(M\) es una variedad lisa orientada de
dimensión cuatro y \(\gamma\in\mathbb R\setminus\{0\}\). Si \(M\) tiene
borde, se imponen variaciones de soporte compacto en el interior o un
término de borde con condiciones que anulen la variación fronteriza. Una
acción de materia espinorial presupone, además, una estructura de espín
compatible.
Si la co-tétrada \(e^I\) toma valores en \(\mathcal L_L\) y la curvatura
\(F^{IJ}\) es una dos-forma sin dimensión después de integrar la conexión,
entonces
\(\int\Sigma_{IJ}\wedge F^{IJ}\) tiene dimensión \(2\mathbf L\).

\begin{theorem}[Homogeneidad de la acción Palatini--Cartan--Holst]
\label{thm:homogeneidad-pch}
Con unidades explícitas, el funcional gravitatorio de dimensión acción es
\[
 \boxed{
 S_{\rm PCH}
 =\frac1{2\kappa_{\rm int}c_{\rm int}}
 \int\Sigma_{IJ}\wedge
 \left(\star F^{IJ}+\gamma^{-1}F^{IJ}\right)
 -\frac1{\kappa_{\rm int}c_{\rm int}}
 \int\Lambda_{\rm int}\operatorname{vol}_e
 +S_{\rm m}.}
 \]
Ambos términos geométricos pertenecen a \(\mathcal L_S\). La escritura con
\(1/(2\kappa)\) sin el factor \(c^{-1}\) pertenece a una convención
\(c=1\) o tiene dimensión \(\mathbf S+\mathbf C\), no dimensión acción.
\end{theorem}

\begin{proof}
De \eqref{eq:dim-Gkappa},
\(\boldsymbol\kappa=-\mathbf S+\mathbf C+2\mathbf T\). Además,
\(2\mathbf L=2\mathbf C+2\mathbf T\). Por tanto,
\[
 -\boldsymbol\kappa-\mathbf C+2\mathbf L
 =\mathbf S.
\]
El término cosmológico posee la misma dimensión porque
\(\boldsymbol\Lambda+4\mathbf L=2\mathbf L\). Si se omite
\(-\mathbf C\), el resultado es \(\mathbf S+\mathbf C\).
\end{proof}

La ecuación de Cartan--Holst conserva entonces su contenido variacional
siempre que la corriente de espín se defina con la normalización coherente
con este funcional. La realización dinámica todavía debe proporcionar un
morfismo desde las holonomías TPK hacia \((e,\omega,\Psi)\); el presente
sección construye el codominio dimensional exacto y elimina la ambigüedad de
las unidades en ese morfismo.

\section{Universalidad de la realización dimensional}

Reunimos los resultados. Sea \(\mathcal X_{\rm HMT}\) el dominio de estados
enriquecidos y sean datos adimensionales sobre él:
\[
 \alpha_{\rm HMT},\quad \Adeg,\quad\Cdeg,\quad
 \widehat c,\widehat Z,\quad
 \mathcal A_*,\quad\chi_{\rm mass},\quad\mathfrak e_0,
\]
construidos en los capítulos anteriores sobre sus dominios declarados.
Consideremos el dominio extendido
\[
 \mathcal X_{\rm MD}
 :=\mathcal X_{\rm HMT}\times\operatorname{Fr}_{\rm D}.
\]

\begin{theorem}[Extensión dimensional universal]
\label{thm:extension-dimensional-universal}
Existe una única extensión tensorial, positiva y
\(\mathcal G_{\rm D}\)-equivariante que:
\begin{enumerate}
 \item realiza el reloj de aristas en \(\mathcal L_T\);
 \item realiza \(\widehat c\) y \(\widehat Z\) en
       \(\mathcal L_C\) y \(\mathcal L_Z\);
 \item realiza el índice determinantal y la componente posterior al retorno
       en \(\mathcal L_S\);
 \item satisface las relaciones exactas entre impedancia, velocidad de
       propagación, permeabilidad y permitividad;
 \item satisface la relación interna de estructura fina con escala elemental
       de carga positiva y, si se proporciona un carácter \(n_Q\), realiza
       condicionalmente \(Q(x)=n_Q(x)e_{\rm int}\);
 \item transporta el carácter de masa en \(\mathcal L_M\).
\end{enumerate}
Todas las demás rectas de la retícula se obtienen por productos tensoriales y
duales. Dos realizaciones difieren exactamente por un único elemento de
\(\mathcal G_{\rm D}\).
\end{theorem}

\begin{proof}
El \cref{thm:base-dimensional-hmt} identifica
\((\mathbf S,\mathbf C,\mathbf T,\mathbf Q)\) con una base libre de la
retícula. Una aplicación monoidal sobre una categoría libre queda determinada
por su valor en los generadores. La positividad fija las raíces que aparecen
en carga y constitutivas. Las identidades ya demostradas garantizan que las
definiciones no dependen de una presentación alternativa de una dimensión.
La acción simplemente transitiva de \(\mathcal G_{\rm D}\) sobre los marcos
prueba la última afirmación.
\end{proof}

La conclusión matemática es precisa. HMT fija la genealogía de los
coeficientes internos y Mecánica Dimensional los realiza en una retícula
ambiental mínima de rango cuatro cuando masa, longitud, duración y carga han
de variar independientemente. El funtor monoidal de trayectorias con firma
ocupa una parte de esa retícula y no hereda, por ello, una afirmación de
minimalidad. Las razones y las ecuaciones homogéneas son independientes de
toda carta. Una publicación en segundos, julios, culombios o
electronvoltios es la composición posterior con un marco externo; no altera
ninguno de los teoremas internos ni puede sustituirlos.
